\section{La crisis ontológica y la arbitrariedad de la métrica nanométrica}

Más allá de su historiografía mitificada, la nanotecnología padece una profunda
crisis ontológica respecto a su definición formal. El estándar internacional, el
cual clasifica como nanotecnología a la manipulación de la materia con al menos
una dimensión entre 1 y 100 nanómetros, es fundamentalmente arbitrario y carece
de justificación física rigurosa. Si bien el límite inferior de 1
nanómetro se corresponde lógicamente con el diámetro cinético de los átomos
ligeros, la cota superior de 100 nanómetros constituye una frontera puramente
administrativa. Esta métrica estrictamente geométrica no
proporciona información alguna sobre la naturaleza física, la complejidad o el
comportamiento dinámico del sistema en cuestión.

Desde la perspectiva de la física fundamental, la transición del régimen
macroscópico, gobernado por la cinemática clásica, a la escala nanométrica,
dominada por la mecánica cuántica y las fluctuaciones térmicas, no ocurre en un
límite invariante de 100 nanómetros. Una definición físicamente
consistente exige como criterio la emergencia de fenómenos cuánticos o
estadísticos que diverjan del comportamiento macroscópico. Este
umbral de transición está dictado por parámetros intrínsecos del material,
tales como el radio de Bohr del excitón, el recorrido libre medio de los
electrones balísticos o la longitud de correlación de un estado fundamental
colectivo. Dependiendo del material y de la fenomenología observada,
estos efectos cuánticos pueden manifestarse a 2 nanómetros o extenderse por
encima de varios cientos de nanómetros.

Al subordinar la naturaleza cuántica de la materia a un umbral geométrico
simplista, la comunidad científica ha facilitado un blanqueamiento
terminológico. En este escenario, procesos químicos ordinarios son
reetiquetados como tecnología de vanguardia por el mero hecho de satisfacer
una restricción espacial.

Esta ausencia de rigor metodológico se exacerba a nivel estatal,
particularmente en ecosistemas científicos en desarrollo. En Colombia, el
intento de establecer un marco regulatorio derivó en una definición
operativa excesivamente laxa. El Consejo Nacional Asesor de Nanociencia y
Nanotecnología propuso clasificar como nanomaterial a cualquier sistema que
exhiba propiedades atribuibles a la escala hasta un límite de un micrómetro.
Adicionalmente, el criterio estipula que un material es
nanoparticulado si apenas el 10 \% de sus partículas presenta un tamaño igual
o inferior a 100 nanómetros, en franco contraste con el umbral del 50 \%
exigido por la Comisión Europea.

La adopción de una métrica que elude tanto el consenso global como las
fronteras estrictas de la física fundamental instituye un andamiaje
burocrático que diluye la integridad del campo. Esta laxitud conceptual
fomenta la apropiación discursiva, permitiendo que la ciencia de materiales
convencional se presente como nanotecnología avanzada sin demostrar un
verdadero dominio atómico.
